\section{Контекстно-свободные языки в качестве ограничений на пути}
	\tikzsetfigurename{FLPQ_CFConstr_}

В данной главе мы дадим обзор задачи достижимости с ограничениями в виде контекстно-свободных языков.
Будут рассмотрены история вопроса, области применения, уточнена формальная постановка для КС-случая и описана общая идея адаптации классических алгоритмов синтаксического анализа к графам.



Задача достижимости с ограничениями в виде контекстно-свободных языков была впервые сформулирована Михалисом Яннакакисом в 1990~году в контексте теории баз данных~\sidecite{Yannakakis}.
Немногим позже Репс с соавторами~\sidecite{10.1145/199448.199462} показал, что широкий класс задач межпроцедурного анализа потока данных может быть точно решён в полиномиальное время путём сведения к достижимости на графах с КС-ограничениями.
Этот подход получил название <<достижимость с ограничениями в виде контекстно-свободных языков>> (\emph{Context-Free Language reachability}, CFL-r) и стал одним из основных инструментов статического анализа программ.
Репс~\sidecite{Reps} предложил общую формулировку межпроцедурного анализа в терминах достижимости на графах.
С тех пор CFL-r активно развивается; современный обзор алгоритической перспективы области представлен Павлогяннисом~\sidecite{10.1145/3583660.3583664}.
Систематическое экспериментальное исследование различных алгоритмов поиска путей с КС-ограничениями (CFPQ) проведено Кёйперсом и~др.~\sidecite{Kuijpers:2019:ESC:3335783.3335791}.

Ключевой прикладной областью CFL-r является межпроцедурный статический анализ кода: задачи анализа формулируются как вопросы достижимости в графе зависимостей программы.
Среди конкретных задач~--- анализ алиасов и указателей~\sidecite{Zheng,10.1145/2001420.2001440,DBLP:conf/asplos/WangHZXS17}, анализ загрязнения~\sidecite{huang2015scalable}, анализ потока меток~\sidecite{LabelFlowCFLReachability}, определение мест исправления ошибок по результатам анализа~\sidecite{CFLfinding} и вывод спецификаций~\sidecite{specificationCFLReachability}.
В биоинформатике КС-достижимость применяется для запросов подграфов в биологических графах~\sidecite{cfpqBio}.
В области семантических веб-технологий и RDF-хранилищ КС-ограничения используются для расширения выразительной силы языков запросов; одним из первых таких расширений является cfSPARQL~\sidecite{10.1007/978-3-319-46523-4_38}.
Нисходящие и матричные методы для КС-запросов к RDF-графам развиваются Медейросом и~др.~\sidecite{MEDEIROS201975,MEDEIROS2022101089}.
Анализ происхождения данных в контексте КС-достижимости рассматривается Миао и Деспанде~\sidecite{MiaoProvenance}.
В системах графовых баз данных, таких как Neo4j~\sidecite{Kuijpers:2019:ESC:3335783.3335791} и RedisGraph~\sidecite{10.1145/3398682.3399163}, ведутся работы по интеграции алгоритмов КС-достижимости в промышленные СУБД.


TODO: (Уточнение формальной постановки для КС-случая) Пусть дан граф $\mscrG = \langle V, E, L \rangle$ и КС-грамматика $G = \langle \Sigma, N, P, S \rangle$, где $L \subseteq \Sigma$.
Задача достижимости: найти множество пар вершин $(v_i, v_j)$, для которых существует путь $\pi$, такой что $\omega(\pi) \in L(G)$.
Альтернативно: для заданных множеств стартовых и финальных вершин.

\subsection{Общая идея адаптации алгоритмов синтаксического анализа к графам}

TODO: Ключевая идея~--- замена позиций в строке на вершины графа.
При обработке терминала в строке есть ровно один следующий символ; в графе~--- несколько исходящих рёбер (конфликт типа shift-shift).
Алгоритмы должны обрабатывать все возможные продолжения.
Различные стратегии обработки этого конфликта приводят к различным алгоритмам КС-достижимости (гл.~\ref{chpt:CFPQ_CYK}, раздел~\ref{sec:GLR}, а также гл.~\ref{chpt:MatrixBasedAlgos}--\ref{chpt:MCFLPQ}).

Проиллюстрируем идею на примере поиска путей в графе с использованием рекурсивного автомата и наивной стратегии вычислений.
\begin{marginfigure}
    \begin{center}
        \resizebox{\marginparwidth}{!}{
                        \begin{tikzpicture}[->,>=stealth']

                \node[initial,state]   (F)              {$q_4$};
                \node[state]           (G) [right = of F] {$q_5$};
                \node[accepting,state] (H) [right = of G] {$q_6$};
                \node[draw=black, fit= (F) (G) (H), inner sep=0.25cm] (J) {};
                \node[below right] at (J.north west) {S'};

                \path (F) edge[below]              node {$q_0$} (G)
                (G) edge[below]              node {$\$$} (H);

                \node[initial,state] (A) [below = 2.6cm of F] {$q_0$};
                \node[state]         (B) [right = of A] {$q_1$};
                \node[state]         (D) [above right = of B] {$q_2$};
                \node[accepting,state]         (C) [right = of B] {$q_3$};
                \node [draw=black, fit= (A) (C) (D), inner sep=0.25cm] (E) {};
                \node[below right] at (E.north west) {S};

                \path (A) edge[below]              node {a} (B)
                (B) edge[left]               node {$q_0$} (D)
                edge[below]              node {b} (C)
                (D) edge[right]              node {b} (C);
            \end{tikzpicture}

        }
    \end{center}
    \caption{Расширенный рекурсивный автомат для грамматики $S \to a \ b \mid a \ S \ b$}
    \label{fig:cfpq-rsm-example}
\end{marginfigure}

\begin{marginfigure}
    \begin{center}
                \begin{tikzpicture}[->,>=stealth',shorten >=1pt,auto,node distance=2.8cm, semithick]

            \node[state] (A)                    {$v_0$};
            \node[state]         (B) [right of=A] {$v_1$};

            \path (A) edge  [loop left] node {a} (A)
            edge  [bend left] node {b} (B)
            (B) edge  [bend left] node {b} (A);
        \end{tikzpicture}

    \end{center}
    \caption{Пример входного графа}
    \label{fig:cfpq-input-graph}
\end{marginfigure}

Пусть на вход подан граф $D$, представленный на рисунке~\ref{fig:cfpq-input-graph}, а грамматика $G$ содержит два правила $S \to a \ b \mid a \ S \ b$.
Расширенный рекурсивный автомат для данной грамматики показан на рисунке~\ref{fig:cfpq-rsm-example}.
Стартовая вершина~--- $v_0$, наша цель~--- найти хотя бы один путь в каждую достижимую (относительно $G$) вершину.

Начальная конфигурация~--- $(q_4,v_0,[])$: мы начинаем из начального состояния блока для $S'$, начальная позиция в графе~--- стартовая вершина $v_0$, стек пуст. На каждом шаге мы применяем правила из определения~\ref{def:rsm_transition} для вычисления новых конфигураций. Последовательность переходов, представленная ниже, позволяет найти путь \[ v_0 \xrightarrow{a} v_0 \xrightarrow{b} v_1 \]
из $v_0$ в $v_1$ (шаг~\ref{eq:cfpq-rsm-step-res-1}) и путь
\[ v_0 \xrightarrow{a} v_0 \xrightarrow{a} v_0 \xrightarrow{b} v_1 \xrightarrow{b} v_0 \]
из $v_0$ в себя (шаг~\ref{eq:cfpq-rsm-step-res-2}).

\begin{align}
    \setcounter{equation}{0}
    (q_4,v_0,[])     \vdash           & \{(q_0,v_0,[q_5])\}                                                   \\
    (q_0,v_0,[q_5])  \vdash           & \{(q_1,v_0,[a,q_5])\}                                                 \\
    (q_1,v_0,[a,q_5])\vdash           & \{(q_0,v_0,[q_2,a,q_5]) \nonumber                                     \\
                                      & , (q_3,v_1,[b,a,q_5])\}                                               \\
    (q_0,v_0,[q_2,a,q_5]) \vdash      & \{(q_1,v_0,[a,q_2,a,q_5])\}                                           \\
    (q_3,v_1,[b,a,q_5])   \vdash      & \boxed{\{(q_5,v_1,[S(a,b)])\}} \label{eq:cfpq-rsm-step-res-1}        \\
    (q_1,v_0,[a,q_2,a,q_5]) \vdash    & \{(q_3,v_1,[b,a,q_2,a,q_5]) \nonumber                                 \\
                                      & , (q_0,v_0,[q_2,a,q_2,a,q_5])\} \label{eq:cfpq-rsm-step-6}           \\
    (q_3,v_1,[b,a,q_2,a,q_5]) \vdash  & \{(q_2,v_1,[S(a,b),a,q_5])\}                                          \\
    (q_2,v_1,[S(a,b),a,q_5]) \vdash   & \{(q_3,v_0,[b,S(a,b),a,q_5])\}                                        \\
    (q_3,v_0,[b,S(a,b),a,q_5]) \vdash & \boxed{\{(q_5,v_0,[S(a,S(a,b),b)])\}} \label{eq:cfpq-rsm-step-res-2}
\end{align}

Заметим, что у нас нет условий для остановки вычислений, из-за чего в нашем примере можно продолжить вычисления и получить новые пути между $v_0$ и $v_1$.
Более того, можно заметить, что таких путей бесконечное количество, а выбор следующего шага не является детерминированным.
Например, можно выбрать конфигурацию $(q_0,v_0,[q_2,a,q_2,a,q_5])$ для продолжения вычислений после шага~\ref{eq:cfpq-rsm-step-6} с риском попасть в бесконечный цикл, но выбрав $(q_3,v_1,[b,a,q_2,a,q_5])$, мы получим новый путь.


\subsection{Представление всех путей: SPPF для графов}

Теоретические основы SPPF изложены в разделе~\ref{sec:SPPF} главы~\ref{chpt:CFG}.
Здесь мы рассмотрим применение SPPF для представления результатов КС-запросов к графам.

Мы, кроме традиционного использования, будем применять SPPF для представления результатов КС запросов к графам.

В графе может существовать множество способов получить путь из одной вершины в другую. И точно так же при построении деревьев вывода путей может появиться несколько одинаковых нетерминалов, получаемых в разных деревьях по-разному. При объединении в SPPF может оказаться, что какой-то путь из вершины $a$ в вершину $b$ является подпутем другого пути из вершины $a$ в вершину $b$, просто более длинного. То есть появятся циклические зависимости.

\begin{example}
    Рассмотрим пример SPPF для задачи поиска путей с КС ограничениями.
    Пусть дан граф $\mathcal{G}:$

    \begin{center}
       \input{figures/edge_labelled_graph_a_3_loop_b_2_loop}
    \end{center}

    Дана грамматика

    \begin{align*}
        G = \langle \{a,b\}, \{S\},  S, \ & \{ \\
            & \ \ (0)\ S  \to a \ S \ b, \\
            & \ \ (1)\ S  \to a \ b, \\
            \ & \} \rangle
    \end{align*}


    Попробуем найти все пути из вершины 2 в вершину 2, выводимые из нетерминала $S$.
    Проверить наличие такого пути можно используя уже известные нам алгоритмы, однако сами пути пока будем строить <<методом пристального взгляда>>.  Найдем один из них. Пусть это будет
    \[ 2 \xrightarrow{a} 0 \xrightarrow{a} 1 \xrightarrow{a} 2 \xrightarrow{a} 0 \xrightarrow{a} 1 \xrightarrow{a} 2 \xrightarrow{b} 3 \xrightarrow{b} 2 \xrightarrow{b} 3 \xrightarrow{b} 2 \xrightarrow{b} 3 \xrightarrow{b} 2.
    \]
    Построим дерево его вывода.

    \begin{center}
    \resizebox{0.3\textwidth}{!}{
        \begin{tikzpicture}[shorten >=1pt,on grid,auto,node distance=1.8cm]
        \node[symbol_node] (s_2_2)   {$(2,S,2)$};

        %\node[prod_node] (p_0_1) [below =of s_0_6] {$0$};
        \node[prod_node] (p_0_1) [below =of s_2_2] {$0$};
        \node[symbol_node] (s_0_3) [below =of p_0_1]   {$(0,S,3)$};
        \node[symbol_node] (a_2_0_1) [left =of s_0_3]   {$(2,a,0)$};
        \node[symbol_node] (b_3_2_1) [right=of s_0_3]   {$(3,b,2)$};

        \node[prod_node] (p_0_2) [below =of s_0_3] {$0$};
        \node[symbol_node] (s_1_2) [below =of p_0_2]   {$(1,S,2)$};
        \node[symbol_node] (a_0_1_1) [left =of s_1_2]   {$(0,a,1)$};
        \node[symbol_node] (b_2_3_1) [right=of s_1_2]   {$(2,b,3)$};

        \node[prod_node] (p_0_3) [below =of s_1_2] {$0$};
        \node[symbol_node] (s_2_3) [below =of p_0_3]   {$(2,S,3)$};
        \node[symbol_node] (a_1_2_1) [left =of s_2_3]   {$(1,a,2)$};
        \node[symbol_node] (b_3_2_2) [right=of s_2_3]   {$(3,b,2)$};

        \node[prod_node] (p_0_4) [below =of s_2_3] {$0$};
        \node[symbol_node] (s_0_2) [below =of p_0_4]   {$(0,S,2)$};
        \node[symbol_node] (a_2_0_2) [left =of s_0_2]   {$(2,a,0)$};
        \node[symbol_node] (b_2_3_2) [right=of s_0_2]   {$(2,b,3)$};

        \node[prod_node] (p_0_5) [below =of s_0_2] {$0$};
        \node[symbol_node] (s_1_3) [below =of p_0_5]   {$(1,S,3)$};
        \node[symbol_node] (a_0_1_2) [left =of s_1_3]   {$(0,a,1)$};
        \node[symbol_node] (b_3_2_3) [right=of s_1_3]   {$(3,b,2)$};

        \node[prod_node] (p_1_1) [below =of s_1_3] {$1$};
        \node[symbol_node] (a_1_2_2) [below left =of p_1_1]   {$(1,a,2)$};
        \node[symbol_node] (b_2_3_3) [below right=of p_1_1]   {$(2,b,3)$};

        \path[->]
        (s_2_2) edge (p_0_1)

        (p_0_1) edge (s_0_3)
        (p_0_1) edge (a_2_0_1)
        (p_0_1) edge (b_3_2_1)

        (s_0_3) edge (p_0_2)

        (p_0_2) edge (s_1_2)
        (p_0_2) edge (a_0_1_1)
        (p_0_2) edge (b_2_3_1)

        (s_1_2) edge (p_0_3)

        (p_0_3) edge (s_2_3)
        (p_0_3) edge (a_1_2_1)
        (p_0_3) edge (b_3_2_2)

        (s_2_3) edge (p_0_4)

        (p_0_4) edge (s_0_2)
        (p_0_4) edge (a_2_0_2)
        (p_0_4) edge (b_2_3_2)

        (s_0_2) edge (p_0_5)

        (p_0_5) edge (s_1_3)
        (p_0_5) edge (a_0_1_2)
        (p_0_5) edge (b_3_2_3)

        (s_1_3) edge (p_1_1)

        (p_1_1) edge (a_1_2_2)
        (p_1_1) edge (b_2_3_3)

        ;
    \end{tikzpicture}

    }
    \end{center}

    Мы построили дерево вывода для одного пути из вершины 2 в неё же.
    Но можно заметить, что таких путей бесконечно много: мы можем бесконечное число раз повторять уже выполненный обход и получать всё более длинные пути.
    В терминах дерева вывода это будет  означать, что к узлу $_1S_3$ мы добавим сына, соответствующего применению продукции 0, а не 1 для нетерминала $S$.
    В таком случае мы получим узел $_2S_2$, который уже существует в дереве и таким образом замкнём цикл.

    \begin{center}
    \resizebox{0.5\textwidth}{!}{
    \input{part_03_GraphAnalysis/chapter_10_FLPQ/figures/05_CFConstraints/05_CFConstraints_fig_02.tex}
    }
    \end{center}

    Таким образом мы построили SPPF. Обойдя эту структуру необходимое количество раз, мы можем получить любой путь, удовлетворяющий условию. Более того, в полученном графе можно получать любые другие пути по соответствующим нетерминалам и парам вершин, содержащимся в узлах леса.
    \end{example}

    %\begin{note}
    SPPF построенный для данной контекстно-свободной грамматики $G$ и графа $\mathcal{G}$
    \begin{enumerate}
      \item содержит терминальный узел вида $(i,t_k,j)$ тогда и только тогда, когда в графе $\mathcal{G}$ есть ребро $(i,t_k,j)$;
      \item содержит нетерминальный узел вида $(i,S_k,j)$ тогда и только тогда, когда в графе $\mathcal{G}$ есть путь из вершины $i$ в вершину $j$, выводимый из нетерминала $S_k$ в грамматике $G$.
    \end{enumerate}
    %\end{note}

    Осталось увидеть, что SPPF является представлением контекстно-свободной грамматики, описывающей результат пересечения исходных графа и грамматики. Для этого просто построим грамматику $G_{\textit{SPPF}} = \langle \Sigma_{\textit{SPPF}}, N_{\textit{SPPF}}, S_{\textit{SPPF}}, P_{\textit{SPPF}}\rangle$ по SPPF следующим образом:
    \begin{itemize}
      \item $\Sigma_{\textit{SPPF}}$~--- все листья SPPF;
      \item $N_{\textit{SPPF}}$~--- все нетерминальные узлы SPPF;
      \item $S_{\textit{SPPF}}$~--- нетерминал, соответствующий пути, который нас будет интересовать;
      \item $P_{\textit{SPPF}}$~--- для каждого дополнительного узла (с номером продукции) добавляем продукцию, левая часть которой~--- непосредственный предок этого узла, а правая чать~--- непосредственные потомки.
    \end{itemize}

    \begin{example}
    Построим грамматику для полученного SPPF:
    \begin{align*}
        (0)\ _2S_2  & \to\ _2a_0\ _0S_3\ _3b_2   &(4)\ _0S_2 \to\ &_0a_1\ _1S_3\ _3b_2 \\
        (1)\ _0S_3  & \to\ _0a_1\ _1S_2\ _2b_3   &(5)\ _1S_3 \to\ &_1a_2\ _2S_2\ _2b_3 \\
        (2)\ _1S_2  & \to\ _1a_2\ _2S_3\ _3b_2   &(6)\ _1S_3 \to\ &_1a_2\ _2b_3 \\
        (3)\ _2S_3  & \to\ _2a_0\ _0S_2\ _2b_3   &
    \end{align*}
    Видим, что для одного единственного нетерминала $_1S_3$ существует 2 правила, одно из которых рекурсивное. Попробуем получить левосторонний вывод какой-нибудь цепочки в этой грамматике:
    \begin{align*}
        & \boldsymbol{_2S_2} \xRightarrow{(0)\ } \\
        & {_2a_0}\ \boldsymbol{_0S_3}\ {_3b_2} \xRightarrow{(1)\ } \\
        & {_2a_0}\ {_0a_1}\ \boldsymbol{_1S_2}\ {_2b_3}\ {_3b_2} \xRightarrow{(2)\ } \\
        & {_2a_0}\ {_0a_1}\ {_1a_2}\ \boldsymbol{_2S_3}\ {_3b_2}\ {_2b_3}\ {_3b_2} \xRightarrow{(3)\ } \\
        & {_2a_0}\ {_0a_1}\ {_1a_2}\ {_2a_0}\ \boldsymbol{_0S_2}\ {_2b_3}\ {_3b_2}\ {_2b_3}\ {_3b_2} \xRightarrow{(4)\ } \\
        & {_2a_0}\ {_0a_1}\ {_1a_2}\ {_2a_0}\ {_0a_1}\ \boldsymbol{_1S_3}\ {_3b_2}\ {_2b_3}\ {_3b_2}\ {_2b_3}\ {_3b_2} \xRightarrow{(5)\ } \\
        & {_2a_0}\ {_0a_1}\ {_1a_2}\ {_2a_0}\ {_0a_1}\ {_1a_2}\ \boldsymbol{_2S_2}\ {_2b_3}\ {_3b_2}\ {_2b_3}\ {_3b_2}\ {_2b_3}\ {_3b_2} \xRightarrow{(0)\ } \\
        & {_2a_0}\ {_0a_1}\ {_1a_2}\ {_2a_0}\ {_0a_1}\ {_1a_2}\ {_2a_0}\ \boldsymbol{_0S_3}\ {_3b_2}\ {_2b_3}\ {_3b_2}\ {_2b_3}\ {_3b_2}\ {_2b_3}\ {_3b_2} \xRightarrow{(1)\ } \\
        & {_2a_0}\ {_0a_1}\ {_1a_2}\ {_2a_0}\ {_0a_1}\ {_1a_2}\ {_2a_0}\ {_0a_1}\ \boldsymbol{_1S_2}\ {_2b_3}\ {_3b_2}\ {_2b_3}\ {_3b_2}\ {_2b_3}\ {_3b_2}\ {_2b_3}\ {_3b_2} \xRightarrow{(2)\ } \\
        & {_2a_0}\ {_0a_1}\ {_1a_2}\ {_2a_0}\ {_0a_1}\ {_1a_2}\ {_2a_0}\ {_0a_1}\ {_1a_2}\ \boldsymbol{_2S_3}\ {_3b_2}\ {_2b_3}\ {_3b_2}\ {_2b_3}\ {_3b_2}\ {_2b_3}\ {_3b_2}\ {_2b_3}\ {_3b_2} \xRightarrow{(3)\ } \\
        & {_2a_0}\ {_0a_1}\ {_1a_2}\ {_2a_0}\ {_0a_1}\ {_1a_2}\ {_2a_0}\ {_0a_1}\ {_1a_2}\ {_2a_0}\ \boldsymbol{_0S_2}\ {_2b_3}\ {_3b_2}\ {_2b_3}\ {_3b_2}\ {_2b_3}\ {_3b_2}\ {_2b_3}\ {_3b_2}\ {_2b_3}\ {_3b_2} \xRightarrow{(4)\ } \\
        & {_2a_0}\ {_0a_1}\ {_1a_2}\ {_2a_0}\ {_0a_1}\ {_1a_2}\ {_2a_0}\ {_0a_1}\ {_1a_2}\ {_2a_0}\ {_0a_1}\ \boldsymbol{_1S_3}\ {_3b_2}\ {_2b_3}\ {_3b_2}\ {_2b_3}\ {_3b_2}\ {_2b_3}\ {_3b_2}\ {_2b_3}\ {_3b_2}\ {_2b_3}\ {_3b_2} \xRightarrow{(6)\ } \\
        & {_2a_0}\ {_0a_1}\ {_1a_2}\ {_2a_0}\ {_0a_1}\ {_1a_2}\ {_2a_0}\ {_0a_1}\ {_1a_2}\ {_2a_0}\ {_0a_1}\ {_1a_2}\ {_2b_3}\ {_3b_2}\ {_2b_3}\ {_3b_2}\ {_2b_3}\ {_3b_2}\ {_2b_3}\ {_3b_2}\ {_2b_3}\ {_3b_2}\ {_2b_3}\ {_3b_2}
    \end{align*}

    Мы получили цепочку, которая действительно является путем из вершины 2 в вершину 2 в заданном графе. Таким образом выводятся и любые другие соответствующие пути.

\end{example}
